\scanpage{288}
\mspage{29}

Mais comme \deleted{\unsure{la}}
\[ \left\{\begin{aligned}\unsure{\sigma|\alpha}(c)&=c(0_{\deleted{\unsure{?}}\tau}),\quad \beta(c)=c(1_\tau)\\
&\text{où }0_\tau,1_\tau\text{ sont les objets source et but}\\
&\text{de l'intervalle orienté \unsure{standard|standart} }I_\tau
\end{aligned}\right. \]
et de $=$ pour $\tau'$, \unsure{$u$|il} \deleted{\unsure{est|n'est}} \unsure{induira|n'induira} bien
\[ \text{\deleted{\unsure{source}}}\quad\left\{\begin{aligned}
\sigma(c')&\longrightarrow c(\varphi(0_{\tau'})\overset{?}{=}\sigma(c)\\[-3pt]
\Vert &\hspace{10em}\Vert\\
c'(0_{\tau'})&\hspace{9em}c(0_\tau)
\end{aligned}\right., \]
et idem pour $\beta(c')$ et $\beta(c)$, que si on impose
\[ \boxed{\left\{\begin{aligned}\varphi(0_{\tau'})&=0_\tau\\ \varphi(1_{\tau'})&=1_\tau\end{aligned}\right.} \]
i.e. que le foncteur envisagé $\varphi$ soit \uline{compatible avec les sources et les buts}. C'est ça \inserted{visiblement} qu'il fallait imposer dans \deleted{\unsure{$\varphi$}} la définition de
\[ \mathrm{Hom}_{\underline{\mathrm{Ty}}}(\tau',\tau)=\underline{\mathrm{Hom}}^{\tikz[baseline=(t.base)]{\node[draw,ellipse,inner sep=.6pt] (t) {\scriptsize int};}}(I_{\tau'},I_\tau)\quad - \]
il faut prendre les foncteurs \inserted{$\varphi$} respectant les structures d'intervalle, i.e. rendant commutatif \unsure{le|les} diagramme
\diagram{\begin{tikzcd}[ampersand replacement=\&,column sep=small]
I_{\tau'} \arrow[rr] \&{} \& I_\tau\\
\deleted{\unsure{\varepsilon}} \&\varepsilon=\{0,1\} \arrow[ul,hook,"{i_{\tau'}}"'] \arrow[ur,hook,"{i_\tau}"']\&{}
\end{tikzcd}}
% slanted
\margin{En fait, il faut sans doute \unsure{restreindre|contraindre} encore plus les morphismes $\varphi$, en exigeant qu'ils soient \uline{croissants} sur les objets --- cf page 31.}

définissant l'orientation de $I_{\tau'}$ et de $I_\tau$. Cette rectification faite, j'ai bon espoir que la catégorie $\underline{\mathrm{Ch}}(X)$ qu'on vient de construire est ``la bonne''. (On pourrait désigner par $\mathrm{CH}(X)$ la catégorie plus grande, sur la catégorie $\underline{\unsure{\mathrm{Typ}|\mathrm{Ty}_p}}$ \deleted{plus} ayant ${=}^{\mathrm s}$ objets que $\underline{\mathrm{Ty}}$, mais plus de morphismes, \unsure{à savoir}
% archive sheet number: 287
